\section{The Mezuro project}
\label{sec:mezuro}

Mezuro is composed of three parts:
(i) the configuration prior to the analysis needless;
(ii) source code metrics computation and evaluation;
(iii) and a graphic interface to present results.
%
Currently, both the computation and visualization modules use other tools
developed into the Mezuro project: Kalibro and Prezento, respectively. In
short, Mezuro is the integration the interaction between these tools, as showed
in Figure \ref{fig:architecture-2}.

Since its first implementation in 2010\cite{mezuro2012}, until being completely
rewritten in 2015, the Mezuro architecture of the system evolved adopting the
microservice architecture\cite{namiot2014micro}, to (i) minimize the amount of
code to maintain; (ii) test and grant quality of code; (iii) modularize the
application in several independent services.

\begin{figure}[hbt]
  \centering
    \includegraphics[width=.8\linewidth]{images/mezuro-architecturev3.png}
  \caption{Current Mezuro architecture.}
  \label{fig:architecture-2}
\end{figure}

The current Mezuro state is specified in Figure \ref{fig:architecture-2}.
Ellipses represents softwares involved and parallelograms the communication
interfaces between them. At Mezuro base we have Kalibro, segmented in three
smaller entities: Kalibro Processor, responsible for processing and evaluating
metrics; Kalibro Configurations, responsible for metrics definitions and
configurations; and Kalibro Client, responsible for interoperate communications
between these entities.  Prezento, the presentation layer, mainly communicates
with Kalibro Processor and Kalibro Configurations, via the Kalibro Client
interface.

\begin{figure}[hbt]
  \centering
    \includegraphics[width=.9\linewidth]{images/new-repository-results.png}
  \caption{Code analysis feature.}
  \label{fig:feature-1}
\end{figure}

Mezuro features can be divided in two groups:

\begin{itemize}
    \item Project (Figure \ref{fig:feature-1})

    \begin{itemize}
        \item Obtain (download and upload) the source code via repositories (Git, Subversion, and Bazaar) or via compressed file;
        \item Code process scheduling (1 day, 2 days, weekly, biweekly, and monthly);
        \item Selection of the metric configuration should be used for each project;
        \item Visualization of measurement results per file via plot graphics with values over-time;
        \item Public results that are accessible by the community.
    \end{itemize}
    \item Configuration
    \begin{itemize}
        \item Metrics configuration cloning and creation;
        \item Statistics about most popular configurations in the community;
        \item Creation of qualitative ranges associated with values of metrics;
	\item Creation of ``reading groups'' to be used in textual interpretation of metrics results;
        \item Combination of native metrics to create composed and more complex analysis.
    \end{itemize}
\end{itemize}

Mezuro has also a strong learning aspect, in which participants are able to
observe other projects, or clone their configurations and definitions, to learn
with them.  This open mutual interaction is interesting to project managers,
software auditors and even an entire developer team. The Mezuro main goal is to
help software developers and maintainers to understand and feel comfortable to
use source code metrics during the development cycles of their projects.

